% !TeX root = ../main.tex
\subsection{تمرین \secnum{۴.۸}، قدرت راه‌اندازی نسبی}
\label{sec:delay-exercise-drive}

\textbf{خواسته.}
قدرت راه‌اندازی یک \lr{Inverter} چهاربرابر و یک
\lr{NAND2} با ظرفیت ورودی \(3C\) را نسبت به
\lr{Inverter} واحد بیابید
\cite[تمرین \secnum{۴.۸}، صفحه ۱۷۷]{weste}.

\textbf{حل.}
نسبت قدرت راه‌اندازی را با \(x\) نمایش می‌دهیم.
ظرفیت ورودی بر حاصل‌ضرب \(g\)
در ظرفیت ورودی \lr{Inverter} واحد تقسیم می‌شود. برای
\lr{Inverter} چهاربرابر، \(g=1\) و \(C_{\mathrm{in}}=12C\)
است. برای \lr{NAND2} با نسبت عرض‌های مرجع، \(g=4/3\)
و ظرفیت هر ورودی در صورت تمرین \(3C\) است. بنابراین داریم.
\[
  x_{\mathrm{INV}}=\frac{12C}{1\cdot3C}=4,
  \qquad
  x_{\mathrm{NAND2}}=\frac{3C}{(4/3)\cdot3C}=\frac34.
\]
پس \lr{NAND2} با وجود داشتن ظرفیت ورودی برابر با مرجع،
سه‌چهارم قدرت راه‌اندازی آن را دارد.
